%------------------------------------------------------------------------------%
\chapter{Plano Cartesiano e Equações}
\label{cap:equacoes}

%------------------------------------------------------------------------------%
\section{Plano Cartesiano}

\begin{teoria}{Plano Cartesiano}

    \begin{minipage}[c]{.6\columnwidth}
        Para identificar e localizar um ponto em um plano podemos usar o sistema de coordenadas Cartesiano.
        \\
        
        Posicionamos duas retas reais perpendiculares uma a outra, chamamos a horizontal de 
        \textbf{eixo~$\mathbf{x}$} e a vertical de \textbf{eixo~$\mathbf{y}$}.
        \\
        
        O ponto onde as retas se cruzam é chamado \textbf{origem}.
        \\
        
        Qualquer ponto do plano é identificado pelo par ordenado $\mathbf{(x,y)}$. 
        \\
        
        O valor $\mathbf{x}$ é chamado \textbf{abscissa} e $\mathbf{y}$ é a \textbf{ordenada}.
    \end{minipage}%        
    \begin{minipage}[c]{.39\columnwidth}
        \addgraphic{.8}{plano-cartesiano}
    \end{minipage}
            
\end{teoria}

%------------------------------------------------------------------------------%
\section{Distância Entre pontos}

\begin{teoria}{Distância entre pontos}
    Para definirmos a distância entre dois pontos $A$ e $B$ conhecendo suas coordenadas $(x_A,y_A)$ e $(x_B,y_B)$,
    usamos o teorema de Pitágoras em um triângulo definido com a hipotenusa ligando os pontos $A$ e $B$ 
    e um cateto horizontal e outro vertical.
    \[ D^2 = (x_B-x_A)^2 + (y_B-y_A)^2  \qquad \text{ou} \qquad D = \sqrt{ (x_B-x_A)^2 + (y_B-y_A)^2 \;} \]
    
    \todo{Incluir figura}
    
    Também podemos usar a notação 
    \[ \dist(A,B) = \sqrt{(x_B-x_A)^2 + (y_B-y_A)^2} \]
    para indicar a distância entre os pontos $A$ e $B$.
\end{teoria}

\begin{exemplos}
    \begin{exemplo}
        Calcule a distância entre os pontos $(4,2)$ e $(7,6)$.
        \[ D^2  = (x_B-x_A)^2 + (y_B-y_A)^2 = (7-4)^2 + (6-2)^2 = 3^2 + 4^2 = 9 + 16 = 25 \]
        \[ D    = \sqrt{25} = 5 \]    
    \end{exemplo}
\end{exemplos}

%------------------------------------------------------------------------------%
\section{Equação da Circunferência}

\begin{observacao}
    Quando resolvemos uma equação queremos todas as soluções possível. Encontrar apenas uma solução em geral não é suficiente.
\end{observacao}

\begin{teoria}{Equação do Círculo ou Circunferência}
    Definimos uma circunferência como o conjunto de \textbf{todos os pontos} que estão a uma distância definida (\textbf{raio}) de um ponto dado (\textbf{centro}).
    \\
    
    Dado um ponto $C=(x_0,\,y_0)$ e uma número positivo $r$ os pontos $P=(x,\,y)$ que pertencem a circunferência de raio $r$ e centro em $C$ obedecem a relação
    \[ \dist(\,P,\,C\,) = r \]
    isso é, dado um ponto $C=(x_0,\,y_0)$ e uma número positivo $r$ a equação da circunferência de raio $r$ e centro em $C$ é
    \[ (x-x_0)^2 + (y-y_0)^2  = r^2 \]
    Todo os pontos $P$ que satisfizerem essa equação estão na circunferência. 
\end{teoria}

%------------------------------------------------------------------------------%
\section{Equação da Reta}

\begin{teoria}{Declividade de uma reta}
    Dados dois pontos $A=(x_a,y_a)$ e $B=(x_b,y_b)$ que pertencem a uma reta $r$ a declividade desta reta é
    \[ m = \frac{\Delta y}{\Delta x} = \frac{y_b-y_a}{x_b-x_a} \]
    Não importa que pontos são escolhidos a declividade da reta é sempre a mesma.
    \\
    
    Se a reta for horizontal \tabto{38mm} $\Delta y = y_b-y_a = 0$ \tabto{69mm} a declividade é nula. \\
    Se a reta for vertical   \tabto{38mm} $\Delta x = x_b-x_a = 0$ \tabto{69mm} a declividade não existe. 
    
\end{teoria}

\begin{teoria}{Equação da reta}
    A equação de uma reta vertical ($m$ não existe) é
    \[ x = a \]
    A equação de uma reta horizontal ($m=0$) é 
    \[ y = b \]
    A equação de uma reta com declividade $m$ que passa pelo ponto $(a,b)$ é
    \[ m = \frac{y-b}{x-a} \qquad \text{ou equivalentemente} \qquad y = b + m(x-a) \]
    
    Escrevendo essa equação de forma resumida temos a equação reduzida da reta
    \[ y = mx + c \]
    
\end{teoria}

\begin{teoria}{Equação geral da reta}
    Qualquer reta pode ser descrita pela equação 
    \[ Ax + By + C = 0 \]
\end{teoria}

%------------------------------------------------------------------------------%
\section{A Reta Tangente}

\begin{teoria}{A Reta Tangente a uma Circunferência}
    A reta tangente a uma circunferência é a reta que toca a circunferência em um único ponto.
\end{teoria}

\begin{teoria}{A Reta Tangente a uma Curva}
    A reta tangente a uma curva qualquer é a reta que ``apenas toca'' a curva em um ponto.
\end{teoria}

%------------------------------------------------------------------------------%
\section{Posição Relativa de Retas}

\begin{teoria}{Posição relativa de retas}
    \begin{description}
        \item[Paralelas]       As retas nunca se encontram
        \item[Concorretes]     As retas se encontram um um ponto
        \item[Perpendiculares] As retas são concorrentes e fazem um ângulo reto ($90^o$) entre elas
    \end{description}    
\end{teoria}

\begin{teoria}{Retas Paralelas}
    Duas retas, $r$ e $s$, são paralelas, se e somente se, suas declividades são iguais 
    \[ m_r = m_s \]
    ou ambas as retas são verticais (a declividade não existe). 
    Se $r$ é paralela a $s$ escrevemos $r\parallel s$.
\end{teoria}

\begin{teoria}{Retas Perpendiculares}
    Duas retas, $r$ e $s$, são perpendiculares, se e somente se, suas declividades obedecem a relação
    \[ m_r = -\frac{1}{\,m_s\,}\]
    ou uma reta é vertical e a outra horizontal.  
    Se $r$ é perpendicular a $s$ escrevemos $r\perp s$.
\end{teoria}

\begin{teoria}{Interseção Entre Retas}
    Duas retas concorrentes tem um ponto em comum, esse ponto satisfaz as equações de ambas as retas.
    Assim, para encontrarmos esse ponto basta resolvermos o sistema de equações lineares
    composto de ambas as equações de reta. Por exemplo se queremos encontrar o ponto de interseção 
    entre as retas $r\!:\, a_r x + b_r y + c_r = 0$ e $s\!:\, a_s x + b_s y + c_s = 0$
    devemos resolver o sistema 
    \[\left\{
        \begin{array}{lclcll}
        a_r x & + & b_r y & + & c_r & = 0 \\
        a_s x & + & b_s y & + & c_s & = 0 
        \end{array}
    \right.
    \]
    Se desejarmos podemos escrever esse sistema na sua forma matricial $AX=B$
    \[
        \left[
        \begin{array}{ll}
            a_r & b_r \\
            a_s & b_s 
        \end{array}
        \right]
        \left[
        \begin{array}{c} x \\ y \end{array}
        \right]
        =
        \left[
        \begin{array}{l} -c_r \\ -c_s \end{array}
        \right]
    \]
    Dessa forma podemos usar o determinante da matriz $A$ para verificarmos se as retas possuem ou não um ponto em comum.
\end{teoria}

%------------------------------------------------------------------------------%
%------------------------------------------------------------------------------%
\section{Exercícios}

\openAnswers{04-Equacoes/04-Equacoes-Exercicios-ans}

\vspace{\baselineskip}
%------------------------------------------------------------------------------%
% \chapter{Plano Cartesiano e Equações}
%------------------------------------------------------------------------------%

%------------------------------------------------------------------------------%
\begin{exercicios}{Plano Cartesiano}
    
    \begin{exercicio}
        Posicione os pontos $(3,2)$, $(4,1)$, $(-1,-1)$, $(\sfrac{1}{2},\sfrac{3}{4})$ e $(\pi,\e)$ nos plano cartesiano.
        \begin{answers} \(\qquad\) \addgraphic{.6}{answer-plano-cartesiano-01}    
        \end{answers}
    \end{exercicio}
    
    \begin{exercicio}
        Escreva os pares ordenados correspondentes aos pontos da figura abaixo.
        \addgraphic{.9}{exercicio-plano-cartesiano}
        \begin{answers}
            \begin{mathlist}[3]
                \item A=(1,8)
                \item B=(-8,8)
                \item C=(-5,6)
                \item D=(-8,3)
                \item E=(-6,-2)
                \item F=(2,-3)
                \item G=(6,-5)
                \item H=(7,4)
                \item I=(5,-2)
                \item J=(4,2)
                \item K=(7,7)
                \item L=(-3,-4)
                \item M=(-7,-6)
                \item N=(7,-2)
                \item P=(7,0)
                \item Q=(-1,0)
                \item R=(0,3)
                \item S=(0,-8)
            \end{mathlist}

        \end{answers}
    \end{exercicio}

    \begin{exercicio}
        Represente, no plano Cartesiano, os pontos 
        $(0,4)$, $(1,0)$, $(1,3)$, $(5,1)$, $(-2,0)$, $(0,-3)$, $(3,-4)$, $(4,-2)$, $(-6,2)$, $(-3,-5)$ e $(-4,-1)$.
        \begin{answers}
            \addgraphic{.6}{answer-plano-cartesiano-03}
        \end{answers}
    \end{exercicio}
    
    \begin{exercicio}
        Se o ponto $P(x,y)$ está no segundo quadrante, quais são os sinais de $x$ e $y$?
    \end{exercicio}
    
    \begin{exercicio}
        Exiba no plano Cartesiano as regiões definidas pelos conjuntos abaixo.
	    \begin{mathlist}[2]
	    	\item \{(x, y) \st x > - 1/2\}
	    	\item \{(x, y) \st -3 \leq y \leq 0\}
	    	\item \{(x, y) \st y \geq 1,5\}
	    	\item \{(x, y) \st x = 3\}
	    	\item \{(x, y) \st -2 \leq x \leq 5\}
	    	\item \{(x, y) \st x \leq -1\}
	    	\item \{(x, y) \st y < 3/2\}
	    	\item \{(x, y) \st y = -2\}
	    \end{mathlist}      
    \end{exercicio}
    
    \begin{exercicio}
        Exiba no plano Cartesiano as regiões definidas abaixo.
	    \begin{mathlist}
	    	\item x \geq 1 \text{ e } y \geq 1
	    	\item x \geq - 1 \text{ e } y \leq 2
	    	\item -3 \leq x \leq 2 \text{ e } 0 \leq y \leq 3
	    	\item -4 \leq x \leq 1 \text{ e } y = 1
	    \end{mathlist}        
    \end{exercicio}
    
    \begin{exercicio}
        Expresse o conjunto de pontos do primeiro quadrante usando desigualdades.
    \end{exercicio}

    \begin{exercicio}
        Verifique, por substituição, se os pontos abaixo pertencem ao gráfico da equação correspondente.
	    \begin{alphalist}
	    	\item \( y = x^3 - 4x^2 + 5x - 12 \)           \tabto{46mm} $(4,10)$                        \tabto{60mm} $(-3,-9)$
	    	\item \( y = \sqrt{x^2 + 1} \)                 \tabto{46mm} $(7,5\sqrt{2})$                 \tabto{60mm} $(0,-1)$
	    	\item \( y = |3x - 8| \)                       \tabto{46mm} $\left(-\sfrac{1}{3},7\right)$  \tabto{60mm} $\left(\sfrac{1}{3}, 7\right)$
	    	\item \( y = \dfrac{1}{|x-2|}-\frac{1}{x-2} \) \tabto{46mm} $\left(-\sfrac{5}{2},0\right)$  \tabto{60mm} $\left(-6,-\!\sfrac{1}{8}\right)$
	    	\item \( y = \dfrac{x^2-2x+1}{-x^3+2x^2+4} \)  \tabto{46mm} $\left(1,\sfrac{2}{3}\right)$   \tabto{60mm} $(2, 4)$
	    	\item \( y^2 + xy^2 - 9 = x^2 + 3xy - 13 \)    \tabto{46mm} $(-2, 6)$                       \tabto{60mm} $(5,-1)$
	    \end{alphalist} 
    \end{exercicio}
    
    \begin{exercicio}
        Usando uma tabela de pares $(x,y)$, trace o gráfico das equações abaixo, no intervalo especificado.
        \begin{mathlist}
            \item y = -2x + 3            \) \tabto{30mm} $x \in [-2, 3]$
            \item 3y - 2x + 3 = 0        \) \tabto{30mm} $x \in [-2, 4]$
            \item 2y + x = 4             \) \tabto{30mm} $x \in [-2, 6]$
            \item y = x^2 - 1            \) \tabto{30mm} $x \in [-2, 2]$
            \item y = 2 - \frac{1}{2}x^2 \) \tabto{30mm} $x \in [-3, 3]$
            \item y = -2x^2 + 4x         \) \tabto{30mm} $x \in [-1, 3]$
            \item y = x^3 - x            \) \tabto{30mm} $x \in [-2, 2]$
            \item y = \sqrt{x}           \) \tabto{30mm} $x \in [ 0, 4]$
            \item y = \sqrt{x + 1}       \) \tabto{30mm} $x \in [-1, 3]$
            \item y = 1/x                \) \tabto{30mm} $x \in [-4, 4]$
            \item y = \abs{x}            \) \tabto{30mm} $x \in [-2, 2]$
            \item y = \abs{x - 2}        \) \tabto{30mm} $x \in [-1, 5]$
        \end{mathlist}   
        \begin{answers}
            \begin{alphalist}[2]
                \item \addgraphic{.6}{answer-plano-cartesiano-09-01}
                \item \addgraphic{.6}{answer-plano-cartesiano-09-02}
                \item \addgraphic{.6}{answer-plano-cartesiano-09-03}
                \item \addgraphic{.6}{answer-plano-cartesiano-09-04}
                \item \addgraphic{.6}{answer-plano-cartesiano-09-05}
                \item \addgraphic{.6}{answer-plano-cartesiano-09-06}
                \item \addgraphic{.6}{answer-plano-cartesiano-09-07}
                \item \addgraphic{.6}{answer-plano-cartesiano-09-08}
                \item \addgraphic{.6}{answer-plano-cartesiano-09-09}
                \item \addgraphic{.6}{answer-plano-cartesiano-09-10}
                \item \addgraphic{.6}{answer-plano-cartesiano-09-11}
                \item \addgraphic{.6}{answer-plano-cartesiano-09-12}
            \end{alphalist}
        \end{answers}
    \end{exercicio}
    
    \begin{exercicio}
        Determine os interceptos e trace o gráfico das equações abaixo.
        \begin{mathlist}[2]
            \item y = x - 1
            \item y = x^2 + 2x - 3
            \item y = 3 - x/2
            \item y = -x^2 + 8x - 12
            \item x = 3 + y
            \item x = y^2 - 1
            \item x = (y+1)(y-2)
            \item x + y = 2
        \end{mathlist} 
        \begin{answers}
            \begin{alphalist}[2]
                \item \addgraphic{.6}{answer-plano-cartesiano-10-01}
                \item \addgraphic{.6}{answer-plano-cartesiano-10-02}
                \item \addgraphic{.6}{answer-plano-cartesiano-10-03}
                \item \addgraphic{.6}{answer-plano-cartesiano-10-04}
                \item \addgraphic{.6}{answer-plano-cartesiano-10-05}
                \item \addgraphic{.6}{answer-plano-cartesiano-10-06}
                \item \addgraphic{.6}{answer-plano-cartesiano-10-07}
                \item \addgraphic{.6}{answer-plano-cartesiano-10-08}
            \end{alphalist}
        \end{answers}
    \end{exercicio}     
    
    \begin{exercicio}
        Represente os pontos em um sistema de eixos coordenados.
        \begin{mathlist}[2]
            \item (-2, 5)
            \item (3, -1)
            \item (3, -4)
            \item (8, -7/2)
            \item (4,5; -4,5)
        \end{mathlist}
        \begin{answers} \(\qquad\) \addgraphic{.6}{answer-plano-cartesiano-11}
        \end{answers}
    \end{exercicio}
   
\end{exercicios}

%------------------------------------------------------------------------------%

%------------------------------------------------------------------------------%
% \chapter{Plano Cartesiano e Equações}
%------------------------------------------------------------------------------%

%------------------------------------------------------------------------------%
\begin{exercicios}{Distância Entre Pontos}
    
    \begin{exercicio}
        Calcule a distância entre os pontos:
        \begin{alphalist}[2]
            \item $(4,1)$ e $(2,3)$
            \item $(1,2)$ e $(-2,-4)$
            \item $(-5,-2)$ e $(3,-2)$
            \item $(-4,5)$ e $(1,-10)$
            \item $(6,4)$ e $(-3,1)$
            \item $(-3,4)$ e $(7,2)$
            \item $(-2,-6)$ e $(-1,1)$
            \item $(5,-2)$ e $(-9,5)$
        \end{alphalist}
        \begin{answers}
            \begin{mathlist}[4]
                \item 2 \sqrt{2}
                \item 3 \sqrt{5}
                \item 8
                \item 5 \sqrt{10}
                \item 3 \sqrt{10}
                \item 2 \sqrt{26}
                \item 5 \sqrt{2}
                \item 7 \sqrt{5}
            \end{mathlist}
        \end{answers}
    \end{exercicio}
    
    \begin{exercicio}
        Qual a distância entre esses pontos e a origem
        \begin{alphalist}[3]
            \item $(4,1)$
            \item $(-2,-4)$
            \item $(-5,-2)$
            \item $(1,-10)$
            \item $(6,4)$
            \item $(-3,4)$
            \item $(-1,1)$
            \item $(5,-2)$
        \end{alphalist}
        \begin{answers}
            \begin{mathlist}[4]
                \item \sqrt{17}
                \item 2 \sqrt{5}
                \item \sqrt{29}
                \item \sqrt{101}
                \item 2 \sqrt{13}
                \item 5
                \item \sqrt{2}
                \item \sqrt{29}
            \end{mathlist}
        \end{answers}
    \end{exercicio}
        
    \begin{exercicio}
        Qual a distância de cada ponto ao ponto~$(-1,1)$
        \begin{alphalist}[3]
            \item $(4,1)$
            \item $(-2,-4)$
            \item $(-5,-2)$
            \item $(1,-10)$
            \item $(6,4)$
            \item $(-3,4)$
            \item $(5,-2)$
        \end{alphalist}
        \begin{answers}
            \begin{mathlist}[4]
                \item 5
                \item \sqrt{26}
                \item 5
                \item 5 \sqrt{5}
                \item \sqrt{58}
                \item \sqrt{13}
                \item 3 \sqrt{5}
            \end{mathlist}
        \end{answers}
    \end{exercicio}
    
    \begin{exercicio}
        Determine a distância entre os pontos dados.
        \begin{alphalist}[2]
            \item $(1, 3)$ e $(4,7)$
            \item $(-1,3)$ e $(4,9)$
        \end{alphalist}	
        \begin{answers}
            \begin{mathlist}[2]
                \item 5
                \item \sqrt{61}
            \end{mathlist}
        \end{answers}
    \end{exercicio}
    
    \begin{exercicio}
        Determine as coordenadas dos pontos  que se encontram a 10 unidades da origem e têm coordenada $y$ igual a $-6$.
    \end{exercicio}
    
    \begin{exercicio}
        Determine as coordenadas dos pontos  que se encontram a 5 unidades da origem e têm coordenada $x$ igual a $3$.
    \end{exercicio}
    
    \begin{exercicio}
        Mostre que os pontos $(3, \, 4)$, $(-3, \, 7)$, $(-6, \, 1)$ e $(0, \, -2)$ formam os vértices de um quadrado.
        \dica{Lembre das propriedades dos lados e das diagonais de um quadrado.}
    \end{exercicio}
    
    \begin{exercicio}
        Mostre que o triângulo com vértices $(-5, \, 2)$, $(-2, \, 5)$ e $(5, \, -2)$ é um triângulo retângulo.
        \dica{Lembre da relação que vale para todo triângulo retângulo.}
    \end{exercicio}
        
    \begin{exercicio}
        Calcule a distância entre os pontos $A(a-3,b+4)$ e $B(a+2,b-8)$.
    \end{exercicio}
    
    \begin{exercicio}
        Calcule o perímetro do triângulo $ABC$, sendo dados $A(2,1)$, $B(-1,3)$ e $C(4,-2)$.
    \end{exercicio}
    
    \begin{exercicio}
        Determine $x$ de modo que o triângulo $ABC$ seja retângulo em $B$, sendo dado que $A(4,5)$, $B(1,1)$ e $C(x,4)$.
    \end{exercicio}
    
    \begin{exercicio}
        Se $P(x,y)$ está a mesma distância de $A(-3,7)$ e $B(4,3)$, qual a relação entre $x$ e $y$?
        \begin{answers}
            \tabto{2em} $14x-8y+33=0$
        \end{answers}
    \end{exercicio}

    \begin{exercicio}
        Dados $A(x,5)$, $B(-2,3)$ e $C(4,1)$ obtenha $x$ de modo que $A$ esteja a mesma distância de $B$ e de $C$.
    \end{exercicio}
    
    \begin{exercicio}
        Dados $(-2,4)$ e $(3,-1)$ vértices consecutivos de um quadrado determine os outros dois vértices.
    \end{exercicio}
    
\end{exercicios}

%------------------------------------------------------------------------------%

%------------------------------------------------------------------------------%
% \chapter{Plano Cartesiano e Equações}
%------------------------------------------------------------------------------%

%------------------------------------------------------------------------------%
\begin{exercicios}{Equação da Circunferência}
    
    \begin{exercicio}
        Determine a equação da circunferência que possui centro em $C(3,6)$ e raio 4.
        \begin{answers}
            \tabto{2em} $(x-3)^2 + (x-6)^2 = 16$
        \end{answers}
    \end{exercicio}
    
    \begin{exercicio}
        Determine a equação da circunferência com raio $2$ e centro em $(-1,3)$ e esboce a circunferência.
    \end{exercicio}

    \begin{exercicio}
        Verifique se os pontos $(6,1)$, $(4,5)$, $(0,5)$, $(5,1)$ e $(-2,-3)$ estão dentro, 
        fora ou sobre a circunferência de raio $5$ e centro $(1,1)$. 
        Esboce a circunferência e os pontos.
    \end{exercicio}
        
    \begin{exercicio}
        Qual a equação da circunferência de raio~$4$ cujo centro é a intersecção das retas 
        \[x -2y +1 = 0 \quad \text{e} \quad x + y - 2 = 0\]
    \end{exercicio}
        
    \begin{exercicio}
        Determine as equações do círculo que satisfazem as condições dadas.
        \begin{alphalist}
            \item Raio $5$ e centro $(2, \, -3)$
            \item Raio $3$ e centro $(-2, \, -4)$
            \item Raio $5$ e centro na origem
            \item Centro $(2, \, -3)$ e passando por $(5, \, 2)$
            \item Centro $(-a, \, a)$ e raio $2a$
        \end{alphalist}
        \begin{answers}
            \begin{mathlist}[2]
                \item (x-2)^2 + (y+3)^2 = 25
                \item (x+2)^2 + (y+4)^2 = 9
                \item x^2 + y^2 = 25
                \item (x-2)^2 + (y+3)^2 = 34
                \item (x+a)^2 + (y-a)^2 = 4a^2 
            \end{mathlist}
        \end{answers}
    \end{exercicio}
    
    \begin{exercicio}
        Determine os coeficientes $a$, $b$ e $c$ da equação do círculo, 
        $ x^2 + y^2 + ax + by + c = 0$, que passa pelos pontos
        $P_1 = (-2,7)$, $P_2 = (-4,5)$ e $P_3 = (4,-3)$.
    \end{exercicio}
    
    \begin{exercicio}
        Construa a equação da circunferência inscrita no quadrado de vértices $(1,1)$, $(3,1)$, $(3,3)$~e~$(1,3)$.
    \end{exercicio}
        
    \begin{exercicio}
        O centro de uma circunferência é determinado pelo ponto médio do segmento 
        $PQ$, sendo $P(4,6)$ e $Q(2,10)$. Considerando que o raio dessa circunferência é 7, 
        determine sua equação.
        \begin{answers}
            \tabto{2em} $(x-3)^2+(y-8)^2=49$
        \end{answers}
    \end{exercicio}
        
    \begin{exercicio}
        O ponto $P(3,b)$ pertence à circunferência de centro no ponto $C(0,3)$ e raio 5. Calcule valor da coordenada $b$. 
        \begin{answers}
            \tabto{2em} $b=-1$ ou $b=7$
        \end{answers}
    \end{exercicio}
    
    \begin{exercicio}
        Determine a equação da circunferência com centro no ponto $C(2,1)$ e que passa pelo ponto $A(1,1)$. 
        \begin{answers}
            \tabto{2em} $(x-2)^2+(y-1)^2=1$
        \end{answers}
    \end{exercicio}

  \begin{exercicio}
      Determine a equação da circunferência de centro $C$ e raio $r$
      \begin{alphalist}
          \item $C(0,0)$                         \tabto{3cm} $r=3$
          \item $C(2,0)$                         \tabto{3cm} $r=4$
          \item $C(-1,-2)$                       \tabto{3cm} $r=5$
          \item $C(2,4)$                         \tabto{3cm} $r=1$
          \item $C(0,-3)$                        \tabto{3cm} $r=2$
          \item $C(\sfrac{1}{2},\,\sfrac{3}{2})$ \tabto{3cm} $r=4$
        \end{alphalist}
        \begin{answers}
            \begin{mathlist}[2]
              \item x ^2 + y^2 = 9 
              \item (x-2) ^2 + y^2 = 16
              \item (x+1) ^2 + (y+2)^2 = 25
              \item (x-2) ^2 + (y-4)^2 = 1 
              \item x ^2 + (y+3)^2 = 4 
              \item (x-\sfrac{1}{2})^2 + (y-\sfrac{3}{2})^2 = 16
            \end{mathlist}
        \end{answers}
    \end{exercicio}
        
    \begin{exercicio}
        Qual é a equação da circunferência de centro $C(1,2)$ que passa por $P(5,5)$?
    \end{exercicio}
    
    \begin{exercicio}
        Determine o centro e o raio das seguintes circunferências
        \begin{mathlist}
            \item (x-2)^2 + (y+1)^2 = 36
            \item (x+2)^2 + y^2 - 9 = 0
            \item y^2 = 16 - x^2
            \item x^2 + y^2 +8x + 6y = 0
            \item x^2 + y^2 -6x + 4y - 12 = 0
        \end{mathlist}
        \begin{answers}
            \begin{mathlist}[2]
                \item C( 2,-1) \qquad r=6
                \item C(-2, 0) \qquad r=3
                \item C( 0, 0) \qquad r=4
                \item C(-4,-3) \qquad r=5
                \item C( 3,-2) \qquad r=5
            \end{mathlist}
        \end{answers}
    \end{exercicio}
        
    \begin{exercicio}
        Esboce as regiões determinadas pelas inequações
        \begin{mathlist}[2]
            \item x^2 + y^2 \leq 9
            \item x^2 + y^2 \geq 4
            \item (x-1)^2 + (1+y)^2 < 9
            \item 1 < (2+x)^2 + y^2 < 4
        \end{mathlist}        
    \end{exercicio}
    
\end{exercicios}

%------------------------------------------------------------------------------%

%------------------------------------------------------------------------------%
% \chapter{Plano Cartesiano e Equações}
%------------------------------------------------------------------------------%

%------------------------------------------------------------------------------%
\begin{exercicios}{Equação da Reta}
    
    \begin{exercicio}
        Encontre a equação da reta que passa pelo ponto $(1,3)$ com declividade~2. 
        Esboce seu gráfico e verifique se os pontos $(2,6)$, $(0,0)$ e $(-1,1)$ pertencem a reta.
    \end{exercicio}

    \begin{exercicio}
        Esboce o gráfico da equação $3x - 4y = 24$.
    \end{exercicio}

    \begin{exercicio}
        Dado que o ponto $(2,-3)$ pertence a reta $-2x+ky+10=0$, determine o valor de $k$.
    \end{exercicio}
    
    \begin{exercicio}
        Encontre as inclinações das retas mostradas na figura abaixo.
        \addgraphic{.6}{exercicio-retas-01}
    \end{exercicio}
        
    \begin{exercicio}
        Determine as inclinações das retas que passam pelos pares de ponto abaixo.
        \begin{alphalist}[2]
            \item $(4,1)$ e $(2,3)$
            \item $(1,2)$ e $(-2,-4)$
            \item $(-5,-2)$ e $(3,-2)$
            \item $(-4,5)$ e $(1,-10)$
            \item $(6,4)$ e $(-3,1)$
            \item $(-3,4)$ e $(7,2)$
            \item $(-2,-6)$ e $(-1,1)$
            \item $(5,-2)$ e $(-9,5)$
        \end{alphalist}
        \begin{answers}
            \begin{mathlist}[4]
                \item m = -1
                \item m = 2
                \item m = 0
                \item m = -3
                \item m = \frac{1}{3}
                \item m = - \frac{1}{5}
                \item m = 7
                \item m = - \frac{1}{2}
            \end{mathlist}
        \end{answers}
    \end{exercicio}
    
    \begin{exercicio}
        Escreva as equações das retas definidas pelas inclinações e interceptos abaixo.
        \begin{alphalist}
            \item Intercepto-$y$:  -1;          inclinação: \sfrac{4}{5}
            \item Intercepto-$y$:   2;          inclinação: \sfrac{-3}{4}
            \item Intercepto-$y$:   4;          inclinação: -3
            \item Intercepto-$y$:  -3;          inclinação: \sfrac{1}{3}
            \item Intercepto-$y$: \sfrac{1}{2}; inclinação:  2
            \item Intercepto-$y$:   0;          inclinação: -1
        \end{alphalist}
        \begin{answers}
            \begin{mathlist}[2]
                \item y = \sfrac{ 4x}{5}    -1
                \item y = \sfrac{-3x}{4}    +2
                \item y = -3x      +4
                \item y = \sfrac{ x}{3}    -3
                \item y =  2x    +\sfrac{1}{2}
                \item y = -x                        
            \end{mathlist}
        \end{answers}
    \end{exercicio}
    
    \begin{exercicio}
        Determine as equações das retas que satisfazem as condições indicadas. Em seguida, trace seus gráficos.
        \begin{alphalist}
            \item Passa por (2,-1)  e tem inclinação 3
            \item Passa por (1,5)   e tem inclinação -3
            \item Passa por (-2,1)  e tem inclinação \sfrac{1}{3}
            \item Passa por (6,-4)  e tem inclinação \sfrac{-1}{3}
            \item Passa por (-4,8)  e tem inclinação -2
            \item Passa por (-3,-2) e tem inclinação \sfrac{3}{2}
        \end{alphalist}
        \begin{answers}
            \begin{mathlist}[3]
                \item y = 3 x - 7
                \item y = - 3 x + 8
                \item y = \sfrac{x}{3} + \sfrac{5}{3}
                \item y = - \sfrac{x}{3} - 2
                \item y = - 2 x
                \item y = \sfrac{3 x}{2} + \sfrac{5}{2}
            \end{mathlist}
        \end{answers}
    \end{exercicio}

    \begin{exercicio}
        Escreva a equação reduzida da reta $r$ que passa pelos pontos $A(-1,1)$ e $B(7,25)$.
    \end{exercicio}
    
    \begin{exercicio}
        Escreva a equação geral das retas que passam pelos pontos
        \begin{alphalist}
            \item $(-2,0)$ \tabto{3cm} $(3,0)$
            \item $(0,-1)$ \tabto{3cm} $(2,0)$
            \item $(-\sfrac{4}{3},0)$ \tabto{3cm} $(0,-2)$
        \end{alphalist}
        \begin{answers}
            \begin{mathlist}[3]
                \item  y + 2  = 0
                \item 2y -  x = 0
                \item 6y + 9x + 20 = 0
            \end{mathlist}
        \end{answers}
    \end{exercicio}
    
    \begin{exercicio}
        Encontre as equações das retas que satisfazem as condições indicadas.
        \begin{alphalist}
            \item Passa por (-1,-3) e intercepta o eixo-$y$ na ordenada 1
            \item Passa por (1,2) e por (2,1)
            \item Passa por (4,-2) e por (-3,-2)
            \item Intercepta o eixo-$y$ na ordenada 3 e o eixo-$x$ na abscissa -2
            \item Intercepta o eixo-$y$ na ordenada 2 e o eixo-$x$ na abscissa 1
            \item Passa por (-2,-6) e intercepta o eixo-$x$ na abscissa 10
            \item Passa por (-6,4) e por (2,-1)
            \item Passa por (-2,8) e por (3,4)
            \item Passa por (-5,0) e por (0,-5)
            \item Passa por (3,14) e por (-4,-28)
            \item Passa por (-1,2 1 ) e por (2,19)
            \item Passa por (3,5;-0,6) e por (-2;2,7)
        \end{alphalist}
    \end{exercicio}
        
    \begin{exercicio}
        Encontre a reta definida pelos pontos 
        \( A\left( \sfrac{7}{2}, \sfrac{5}{2}\right) \) e  
        \( B\left(-\sfrac{5}{2},-\sfrac{7}{7}\right) \).
    \end{exercicio}
        
    \begin{exercicio}
        A reta determinada por $A(a,0)$ e $B(0,b)$ passa por $C(3,4)$. Qual a relação entre $a$ e $b$?
    \end{exercicio}
        
    \begin{exercicio}
        A reta determinada por $A(q,q)$ e $B(3,-2)$ passa pela origem. Qual a relação entre $p$ e $q$?
    \end{exercicio}
    
    \begin{exercicio}
        Desenhe no plano cartesiano as retas cujas equações são
        \begin{mathlist}[2]
            \item y=2x
            \item x+y+3=0
            \item x+y=5
            \item 2y+x=0
            \item x-y+5=0
            \item x-y-4=0
        \end{mathlist}
    \end{exercicio}
    
    \begin{exercicio}
        Determine as equações das retas mostradas na figura abaixo.
        \addgraphic{.6}{exercicio-retas-02}
    \end{exercicio}
    
    \begin{exercicio}
        Trace os gráficos das equações abaixo.
	    \begin{mathlist}[2]
	    	\item y = - \frac{2}{3}x + 1
	    	\item y = 5x - 2
	    	\item y = -2x
	    	\item y = 4
	    	\item x - y = -3
	    	\item 3y - x + 4 = 0
	    \end{mathlist}        
    \end{exercicio}
    
    \begin{exercicio}
        Encontre as equações das retas que satisfazem as condições abaixo. Em seguida, trace os gráficos correspondentes.
        \begin{alphalist}
            \item Passa por (-2,1) e (-2, 5).
            \item Passa por (-7,8) e (10, 8).
            \item Reta vertical que passa por (3,-1).
            \item Reta horizontal que passa por (6,-4).
        \end{alphalist}
    \end{exercicio}
    
    \begin{exercicio}
        Dados os pontos $A(-1,2)$ e $B(3,-1)$
        \begin{alphalist}
            \item Marque os pontos no plano Cartesiano, considerando as abscissas no intervalo $[-3,5]$ e as ordenadas em $[-2, 3]$.
            \item Determine a equação da reta que passa pelos pontos. Trace essa reta no gráfico.
            \item Determine a ordenada do ponto dessa reta no qual a abscissa vale 1.
            \item Determine a abscissa do ponto da reta que tem ordenada 0.
        \end{alphalist}
    \end{exercicio}
    
\end{exercicios}

%------------------------------------------------------------------------------%

%------------------------------------------------------------------------------%
% \chapter{Plano Cartesiano e Equações}
%------------------------------------------------------------------------------%

%------------------------------------------------------------------------------%
\begin{exercicios}{Posição Relativa de Retas}
    
    \begin{exercicio}
        Se a reta que passa pelos pontos $(1,a)$ e $(4,-2)$ é paralela a
        reta que passa pelos pontos $(2,8)$ e $(-7,a+4)$ qual o valor de $a$?
    \end{exercicio}
    
    \begin{exercicio}
        Determinar o ponto onde as retas $3y+3x-2=0$ e $3y-x=1$ se cruzam.
    \end{exercicio}
    
    \begin{exercicio}
        Encontre uma equação para a reta que passa pelo ponto $(2, 3)$ e é  paralela à reta com equação $3x + 4y - 8 = 0$. 		
    \end{exercicio}
    
    \begin{exercicio}
        Encontre uma equação para a reta que passa pelo ponto $(-1, 3)$ e é paralela à reta que passa pelos pontos $(-3, 4)$ e $(2, 1)$.	
    \end{exercicio}

    \begin{exercicio}
        Encontre uma equação para a reta que passa pelo ponto $(-2,-4)$ e é perpendicular à reta com equação $2x-3y-24=0$. 
    \end{exercicio}

    \begin{exercicio}
        Encontre uma equação para a reta que passa pelo ponto $(-2,-4)$ e é perpendicular à reta com equação $2x-3y-24=0$.
    \end{exercicio}
    
    \begin{exercicio}
        Construa a equação da reta $s$ que passa pelo ponto $(1,-1)$ e é perpendicular a reta $r$, definida pela equação $2y=x+4$.
    \end{exercicio}

    \begin{exercicio}
        Determine a interseção entre as retas $x+2y=3$ e $2x+3y=5$.
    \end{exercicio}
    
    \begin{exercicio}
        Mostrar que as retas $2x+3y-1=0$, $x+y=0$ e $3x+4y-1=0$ se encontram no mesmo ponto.
    \end{exercicio}
    
\end{exercicios}

%------------------------------------------------------------------------------%

%------------------------------------------------------------------------------%
% \chapter{Plano Cartesiano e Equações}
%------------------------------------------------------------------------------%

%------------------------------------------------------------------------------%
\begin{exercicios}{Revisão}
    
    \begin{exercicio}
        Qual a equação da circunferência de raio~$3$ cujo centro é a intersecção das retas $x -2y +1 = 0$ e $x + y - 2 = 0$.
    \end{exercicio}
        
    \begin{exercicio}
        Construa a equação da reta $s$ que passa pelo ponto $(2,-2)$ e é perpendicular a reta $r$, definida pela equação $2y=x+4$. 
    \end{exercicio}
        
    \begin{exercicio}
        Dada a circunferência $c$ definida pela equação $(x-1)^2 + (y+1)^2=9$, construa as retas $r$ e $s$ tangentes
        a circunferência nos pontos $(4,-1)$ e $(1,2)$, respectivamente. Qual o ponto de intersecção das retas $r$ e $s$?
        (Dica: esboce a circunferência e as retas.)
    \end{exercicio}
        
    \begin{exercicio}
        Dada a circunferência $c$ definida pela equação $(x-1)^2 + (y+1)^2=4$, construa as retas $r$ e $s$ tangentes
        a circunferência nos pontos $(-1,-1)$ e $(1,3)$, respectivamente. Qual o ponto de intersecção das retas $r$ e $s$?
    \end{exercicio}
        
    \begin{exercicio}
        Dada a circunferência $c$ definida pela equação $(x-1)^2 + (y+1)^2=4$, construa as retas $r$ e $s$ tangentes
        a circunferência nos pontos $(-1,-1)$ e $(1,-3)$, respectivamente. Qual a distância do ponto de intersecção 
        das retas $r$ e $s$ até o centro da circunferência?
    \end{exercicio}
        
    \begin{exercicio}
        Sejam
        \begin{enumerate}[label={\roman*.}]
            \item a circunferência $u$ dada por \[(x-1)^2+(y-2)^2=4\]
            \item a circunferência $v$ dada por \[(x+2)^2=9-y^2\]
            \item a reta $r$ que passa pelos pontos \[(0,-3) \quad \text{ e } (3,-1)\]
            \item a reta $s$ dada pela equação \[2y+3x=5\]
        \end{enumerate}
        \begin{alphalist}
            \item Ilustre as formas geométricas definidas.
            \item Calcule a distância entre os centros das circunferências $u$ e $v$.
            \item Escreva a equação da reta $a$ que passa pelos centros das circunferências $u$ e $v$.
            \item Qual a declividade da reta $a$?
            \item Qual a equação da reta $r$?
            \item Qual a relação entre a reta $r$ e a reta $a$ (elas são paralelas, concorrentes, perpendiculares)?
            \item Escreva a equação da circunferência com centro na origem e que passe pelo centro da circunferência $u$.
            \item Escreva a equação da reta paralela a reta $r$ que passa pela origem.
        \end{alphalist}
    \end{exercicio}
    
    \begin{exercicio}
        Encontre uma equação para a reta $L$ que passa pelo ponto $(-2, 4)$ e satisfaz a condição dada.
        \begin{alphalist}
            \item $L$ é uma reta vertical.
            \item $L$ é uma reta horizontal.
            \item $L$ passa pelo ponto $\left(3, \sfrac{7}{2}\right)$.
            \item A intersecção $x$ de $L$ é $3$.
            \item $L$ é paralela à reta $5x - 2y = 6$.
            \item $L$ é perpendicular à reta $4x + 3y = 6$.
        \end{alphalist}
    \end{exercicio}
  
    \begin{exercicio}
        Sabendo que formula para calcular a distância entre uma reta, $r:ax+by+c=0$, e um ponto, $P(x_0,y_0)$, é
        \[d(r,P) = \frac{|ax_0+by_0+c|}{\sqrt{a^2+b^2}}\]
        Calcule a distância entre a reta e o ponto dados
        \begin{mathlist}
            \item r: 3x+4y+2=0 \) \tabto{33mm} P(-1,1)
            \item r: x+y-1=0   \) \tabto{33mm} P(2,1)
            \item r: y=2x+1    \) \tabto{33mm} P(0,-1)
            \item r: y=x       \) \tabto{33mm} P(0,1)
        \end{mathlist}
    \end{exercicio}

    \begin{exercicio}
        Sabemos que a reta tangente a uma circunferência é perpendicular ao seu raio.
        Usando essa informação, calcule a equação da reta tangente a circunferência 
        $x^2+y^2-25=0$ que passa pelo ponto $(3,4)$.
    \end{exercicio}
        
    \begin{exercicio}
        Determine se as afirmações dadas são verdadeiras ou falsas. 
        Se verdadeiras, explique porquê. Se falsas, justifique através de um exemplo.
        \begin{alphalist}
            \item O ponto $(-a, \, -b)$ é simétrico ao ponto $(a, \, b)$ em relação à origem.
            \item Se a distância entre os pontos $P_1 (a, \, b)$ e $P_ 2 (c, \, d)$ é $D$, 
                então a distância entre os pontos 
                $P_1 (a, \, b)$ e $P_3 (kc, \, kd)$, $(k \, \neq \, 0)$, é dada por $|k| D$.
            \item O círculo com equação $kx^2 \, + \, ky^2 \, = \, a^2$ situa-se dentro do 
                círculo com equação $x^2 \, + \, y^2 \, = \, a^2$, desde que $k \, > \, 1$.
            \item Sejam $(x_1, y_1)$ e $(x_2, y_2)$ dois pontos no plano $xy$. 
                Mostre que a distância entre os dois pontos é dada por
                \[d = \sqrt{(x_2 - x_1)^2 + (y_2 - y_1)^2}\]
        \end{alphalist}        
    \end{exercicio}
        
    \begin{exercicio}
        Determine o perímetro do triângulo $ABC$ que é descrito as seguintes condições
        \begin{alphalist}
            \item O vértice $A$ pertence ao eixo-$x$
            \item O vértice $B$ pertence ao eixo-$y$
            \item A reta que contêm $BC$ tem equação $x-y=0$
            \item A reta que contêm $AC$ tem equação $x+2y-3=0$
        \end{alphalist}
        \begin{answers}
            \tabto{2em} $3+\sqrt{2}+\sqrt{5}$
        \end{answers}
    \end{exercicio}
        
    \begin{exercicio}
        Determine a posição relativa de cada par de retas
        \begin{alphalist}
            \item $2x- y+3=0$  \tabto{3cm}  $4x-2y=-6$
            \item $2x- y+5=0$  \tabto{3cm}  $3x-6y=-3$
            \item $ x-2y+3=0$  \tabto{3cm}  $2x+4y+ 3=0$
            \item $2x- y+3=0$  \tabto{3cm}  $4x-2y+10=0$
            \item $2x+4y+3=0$  \tabto{3cm}  $4x-2y=-6$
        \end{alphalist}
    \end{exercicio}
        
    \begin{exercicio}
        Para quais valores de $k$ as retas $(k-1)x+6y+1=0$ e $4x+(k+1)y-1=0$ são paralelas?
    \end{exercicio}
    
    \begin{exercicio}
        Ache a equação da reta que passa pelo centro da circunferência $(x-3)^2+(y-2)^2=8$ e é perpendicular à reta $x-y-16=0$.
    \end{exercicio}
    
    \begin{exercicio}
        Um quadrado tem vértices consecutivos $A(5,0)$ e $B(-1,0)$. Determine as equações das circunferências inscrita e circunscrita ao quadrado.
    \end{exercicio}
    
    \begin{exercicio}
        Encontre a região do plano cujos pontos $P(x,y)$ satisfazem as relações $x+y\leq 2$ e $x^2+y^2\leq 16$, esboce essa região.
    \end{exercicio}
    
    \begin{exercicio}
        Resolva os sistemas de inequações
        \begin{mathlist}[1]
            \item \systeme{ x^2 + y^2 \leq 25, x^2 + y^2 \geq  4 }
            \item \systeme{ x^2 + y^2 \leq  9, x\phantom{^2} + y\phantom{^2} \leq 3 }
            \item \systeme{ x^2 + y^2 \geq  1, x^2 + y^2 \leq  9 }
            \item \systeme{ x^2 + y^2 \leq  1, x\phantom{^2} + y\phantom{^2} \geq 1 }
            \item \begin{cases}
                   x^2 + y^2     &\!\!\! \leq  4 \\
                   (x-1)^2 + y^2 &\!\!\! \leq  4
                  \end{cases}
            \item \begin{cases}
                x^2 + y^2         &\!\!\! \leq 25 \\
                (x-1)^2 + (y-1)^2 &\!\!\! \geq 16
            \end{cases}
        \end{mathlist}
    \end{exercicio}
    
    \begin{exercicio}
        Sejam os conjuntos 
        \begin{align*}
            A &= \{ (x,y) \st x^2+y^2 \leq 25 \} \\
            B &= \{ (x,y) \st (x-1)^2 + (y-1)^2 \leq 16 \} 
        \end{align*}
        Determine 
        \begin{mathlist}[2]
            \item A \cup B
            \item A \cap B
            \item A - B
            \item B - A
        \end{mathlist}
    \end{exercicio}
        
    \begin{exercicio}
        Encontre os pontos onde a circunferência $(x-1)^2+(y-2)^2=29$ encontra o eixo-$x$.
    \end{exercicio}
    
    \begin{exercicio}
        Encontre os pontos onde a circunferência $(x-1)^2+(y-2)^2=26$ encontra o eixo-$y$.
    \end{exercicio}
    
\end{exercicios}

%------------------------------------------------------------------------------%
\begin{exercicios}{Desafio}
    
    \begin{exercicio}
        Se $r$ e $s$ são retas paralelas e $AX=B$ é o sistema linear construído para 
        calcular o ponto de interseção entre elas, determine $\det A$.
    \end{exercicio}
    
\end{exercicios}
%------------------------------------------------------------------------------%



\writeAnswers{04-Equacoes/04-Equacoes-Exercicios-ans}

%------------------------------------------------------------------------------%
